\appendix
\clearpage
\phantomsection
\refstepcounter{section}
\addcontentsline{toc}{section}{APÊNDICE \Alph{section} -- SÍNTESE DA AVALIAÇÃO AUTOMATIZADA}
\begin{center}
\textbf{APÊNDICE \Alph{section} -- SÍNTESE DA AVALIAÇÃO AUTOMATIZADA}
\end{center}
\label{sec:apendice-avaliacao}

Este apêndice reúne a síntese das três rodadas automatizadas utilizadas para complementar a demonstração empírica apresentada na seção de experimentos, com o objetivo de tornar visíveis, em formato legível, os resultados quantitativos após o alinhamento entre a automação e a configuração operacional do Open WebUI. As categorias apresentadas na tabela combinam tipos de consulta inspirados em benchmarks de avaliação de RAG, especialmente o RAGEval \parencite{zhuEtAl2025rageval}, com extensões metodológicas definidas no próprio experimento para capturar riscos típicos do domínio parlamentar.

\subsection{Síntese das rodadas automatizadas}

A Tabela~\ref{tab:rodadas-avaliacao} apresenta a síntese das três rodadas principais consideradas na análise. Em todas elas, 20 perguntas foram avaliadas com estratégia de \textit{LLM as a Judge} \parencite{liuEtAl2025judgeRag}, mantendo a mesma base de conhecimento, o mesmo conjunto de perguntas e o mesmo modelo nas etapas de geração e avaliação. A nota total de cada item corresponde à soma dos cinco critérios da rubrica, em escala de 0 a 10 \parencite{liuEtAl2023geval}.

\begingroup
\footnotesize
\begin{longtable}{@{}>{\raggedright\arraybackslash}p{0.26\textwidth}cccccc@{}}
\caption{Síntese das três rodadas principais de avaliação automatizada}
\label{tab:rodadas-avaliacao}\\
\textbf{Rodada} & \textbf{Perguntas} & \textbf{Falhas} & \textbf{Média} & \textbf{Mín.} & \textbf{8--10} & \textbf{5--7} \\
\hline
\endfirsthead
\textbf{Rodada} & \textbf{Perguntas} & \textbf{Falhas} & \textbf{Média} & \textbf{Mín.} & \textbf{8--10} & \textbf{5--7} \\
\hline
\endhead
Primeira & 20 & 0 & 9,15 & 6 & 17 & 3 \\
Segunda & 20 & 0 & 9,55 & 7 & 19 & 1 \\
Terceira & 20 & 0 & 9,35 & 7 & 19 & 1 \\
\hline
\end{longtable}
\endgroup

A Tabela~\ref{tab:resultado-final-20} detalha as notas obtidas por pergunta nas três rodadas. Os resultados completos estão disponíveis nos arquivos públicos da primeira\footnote{\url{https://github.com/fabriciosantana/mcdia/blob/main/05-iag/4-project/eval/results/rag_eval_20260416T172816Z.csv}}, segunda\footnote{\url{https://github.com/fabriciosantana/mcdia/blob/main/05-iag/4-project/eval/results/rag_eval_20260416T182240Z.csv}} e terceira\footnote{\url{https://github.com/fabriciosantana/mcdia/blob/main/05-iag/4-project/eval/results/rag_eval_20260416T232921Z.csv}} rodadas.

\begingroup
\footnotesize
\setlength{\LTpre}{0pt}
\setlength{\LTpost}{0pt}
\begin{longtable}{@{}p{0.07\textwidth}>{\raggedright\arraybackslash}p{0.15\textwidth}>{\raggedright\arraybackslash}p{0.52\textwidth}ccc@{}}
\caption{Notas por pergunta nas três rodadas automatizadas principais}\label{tab:resultado-final-20}\\
\textbf{ID} & \textbf{Categoria} & \textbf{Pergunta} & \textbf{R1} & \textbf{R2} & \textbf{R3} \\
\hline
\endfirsthead
\textbf{ID} & \textbf{Categoria} & \textbf{Pergunta} & \textbf{R1} & \textbf{R2} & \textbf{R3} \\
\hline
\endhead
q01 & factual & Quais dados Ciro Nogueira apresenta para criticar a concentração bancária no Brasil e qual conclusão ele tira sobre juros, tarifas e spread? & 10 & 10 & 9 \\
q02 & factual & Quais são os principais argumentos de Paulo Paim contra a capitalização da Previdência e que exemplos ele usa para sustentar essa posição? & 10 & 9 & 9 \\
q03 & proposta & Que propostas Confúcio Moura apresenta para melhorar a educação? & 10 & 10 & 10 \\
q04 & proposta & Por que Confúcio Moura defende que senadores e deputados acompanhem mais de perto escolas públicas? & 10 & 10 & 10 \\
q05 & foco no autor & Quais prioridades Chico Rodrigues afirma defender no início do mandato e como ele relaciona ordem, união e desenvolvimento para o Brasil e para Roraima? & 10 & 10 & 10 \\
q06 & comparação & Compare as preocupações de Paulo Paim sobre Previdência com as preocupações de Confúcio Moura sobre educação. & 10 & 10 & 10 \\
q07 & síntese & Quais temas sociais e econômicos aparecem com mais frequência nos discursos carregados e quais senadores associam esses problemas a propostas concretas de solução? & 9 & 10 & 9 \\
q08 & números & Quais discursos citam percentuais, valores ou estatísticas e para sustentar quais argumentos? & 8 & 10 & 10 \\
q09 & foco no autor & Que argumentos Wellington Fagundes usa para defender investimentos em infraestrutura e qual relação ele faz entre logística e competitividade? & 10 & 10 & 8 \\
q10 & ampla & Quais diagnósticos sobre desigualdade, vulnerabilidade ou falta de desenvolvimento aparecem nesses discursos e que tipos de solução são defendidos? & 10 & 9 & 9 \\
q11 & comparação & Compare como Ciro Nogueira e Wellington Fagundes usam dados econômicos ou setoriais para defender mudanças estruturais. Quais números aparecem, para sustentar que tese, e onde os argumentos deles divergem? & 8 & 9 & 9 \\
q12 & desambiguação & Quais propostas sobre educação aparecem especificamente em falas de Confúcio Moura e quais aparecem em falas de outros senadores, mas poderiam ser confundidas com as dele? Responda separando com clareza o que é dele e o que não é. & 10 & 10 & 10 \\
q13 & evidência negativa & Há, no contexto recuperado, evidência suficiente para afirmar que existe uma proposta detalhada de implementação, com prazos e fonte de financiamento, para combater a desigualdade? Se não houver, explique exatamente o que falta. & 10 & 10 & 10 \\
q14 & multietapas & A partir de diferentes discursos, reconstrua a cadeia de raciocínio que liga desigualdade, desenvolvimento regional, educação e proteção social. Quais autores fazem essas conexões de modo mais explícito e com quais limites? & 7 & 8 & 7 \\
q15 & precisão numérica & Quais números citados nos discursos parecem centrais para o argumento de cada autor e quais aparecem apenas como apoio periférico? Distingua os dois grupos e evite misturar estatísticas de autores diferentes. & 6 & 10 & 10 \\
q16 & controle de escopo & O que exatamente Paulo Paim afirma sobre capitalização da Previdência, e o que seria uma extrapolação indevida para além do que os trechos permitem concluir? & 7 & 9 & 10 \\
q17 & foco na fonte & Quais respostas exigiriam mais cautela por dependerem de inferências amplas a partir de poucos trechos: as sobre educação, as sobre desigualdade ou as sobre infraestrutura? Justifique com base nas limitações do próprio contexto. & 10 & 10 & 10 \\
q18 & autoria cruzada & Quais semelhanças de linguagem ou tema entre Paulo Paim, Confúcio Moura e outros autores podem levar o sistema a confundir autoria ou proposta? Identifique os riscos concretos de confusão e ilustre com exemplos. & 8 & 7 & 9 \\
q19 & estresse de recuperação & Se a pergunta for ``quais senadores apresentam diagnóstico e solução no mesmo discurso, sem depender de síntese entre vários discursos?'', quem realmente se encaixa nisso no contexto disponível? Cite apenas os casos sustentados com segurança. & 10 & 10 & 8 \\
q20 & checagem de alucinação & Entre as perguntas deste conjunto, quais tenderiam a induzir alucinação ou excesso de confiança se o modelo não reconhecesse limites do contexto? Explique por que e indique como uma resposta cautelosa deveria se comportar. & 10 & 10 & 10 \\
\hline
\end{longtable}
\endgroup

\subsection{Rubrica utilizada na avaliação automatizada}

Para favorecer reprodutibilidade, a seguir registra-se o conteúdo da rubrica aplicada pelo LLM juiz. Sua formulação em critérios explícitos e escala discreta aproxima o procedimento adotado de abordagens de avaliação por formulário e rubrica descritas na literatura de \textit{LLM-based evaluation} \parencite{liuEtAl2023geval}. As faixas de decisão registradas na rubrica têm sentido operacional no escopo do experimento, não devendo ser lidas como certificação de prontidão para uso institucional amplo.

\begin{Verbatim}[breaklines=true]
# RAG Evaluation Rubric

Use esta rubric para classificar cada resposta em uma planilha ou revisao manual.

## Escala

- `2`: bom
- `1`: aceitavel
- `0`: ruim

## Criterios

### 1. Aderencia a pergunta
- `2`: responde exatamente ao que foi perguntado
- `1`: responde parcialmente ou abre demais o escopo
- `0`: foge do foco

### 2. Precisao factual
- `2`: fatos, numeros, autores e temas corretos
- `1`: pequenos desvios ou formulacoes muito amplas
- `0`: erros factuais relevantes

### 3. Foco nas fontes
- `2`: usa trechos coerentes e pertinentes
- `1`: usa fontes validas, mas mistura demais ou inclui evidencias perifericas
- `0`: cita fontes incompatíveis ou aparentemente sem relacao

### 4. Sintese
- `2`: bem organizada, clara e proporcional
- `1`: correta, mas excessivamente longa, repetitiva ou burocratica
- `0`: confusa ou mal estruturada

### 5. Confianca / Alucinacao
- `2`: nao inventa e sinaliza limites quando necessario
- `1`: faz inferencias amplas, mas sem erros graves
- `0`: aparenta inventar ou extrapolar sem base

## Decisao pratica

- `8 a 10`: pronta para uso
- `5 a 7`: utilizavel com supervisao
- `0 a 4`: precisa ajuste de prompt, retrieval ou pergunta

## Padroes comuns de erro

- Mistura excessiva de autores ou discursos
- Resposta mais ampla que a pergunta
- Generalizacao sem apoio suficiente
- Falta de citacoes ou citacoes mal distribuídas
- Nao reconhecer quando faltou informacao no contexto
\end{Verbatim}
